% Source: https://artofproblemsolving.com/community/c3h3134466_2012_lmt_spring_guts_round_p13p24_lexington_mathematical_tournament
% Stable identity: aops-topic-3134466-q15
The equal-side pairs make $BD$ an axis of symmetry, so $\angle A=\angle C$. Cyclicity makes them supplementary, hence both right angles. Thus $BD=\sqrt{1^2+2^2}=\sqrt5$ is a circumdiameter, and $R^2=\dfrac{5}{4}$. The quadrilateral consists of two right triangles of area 1 each, so its area is 2. Its semiperimeter is 3, so the inradius is $r=\dfrac{2}{3}$ by the sum of four base-times-radius triangle areas. The requested area is $\pi(R^2-r^2)=\pi(\dfrac{5}{4}-\dfrac{4}{9})=\boxed{\dfrac{29\pi}{36}}$.
\begin{asy}
usepackage("amsmath");
size(7cm);
pair B=(0,0),D=(sqrt(5),0),A=(1/sqrt(5),2/sqrt(5)),C=(1/sqrt(5),-2/sqrt(5));draw(circle((sqrt(5)/2,0),sqrt(5)/2));draw(A--B--C--D--cycle);draw(circle((sqrt(5)/3,0),2.0/3),blue);label("$A$",A,N);label("$B$",B,W);label("$C$",C,S);label("$D$",D,E);
\end{asy}
